\begin{thebibliography}{99}
\bibitem[1]{source1} Martín Abadi and Véronique Cortier. Deciding knowledge in security protocols under equational theories. Theoretical Computer Science, 367(1-2):2-32, 2006. URL: https://doi.org/10.1016/J.TCS.2006.08.032 .
\bibitem[2]{source2} Martín Abadi and Cédric Fournet. Mobile values, new names, and secure communication. In Chris Hankin and Dave Schmidt, editors, Conference Record of POPL 2001: The 28th ACM SIGPLAN-SIGACT Symposium on Principles of Programming Languages, London, UK, January 17-19, 2001, pages 104-115. ACM, 2001. URL: https://doi.org/10.1145/360204.360213 .
\bibitem[3]{source3} Zümrüt Akçam, Kimberly A Cornell, Daniel S Hono II, Paliath Narendran, and Andrew Pulver. On problems dual to unification: The string-rewriting case. arXiv preprint, 2021. URL: https://arxiv.org/abs/2103.00386 .
\bibitem[4]{source4} Siva Anantharaman, Hai Lin, Christopher Lynch, Paliath Narendran, and Michaël Rusinowitch. Unification modulo homomorphic encryption. Journal of automated reasoning, 48(2):135-158, 2012. URL: https://doi.org/10.1007/S10817-010-9205-Y .
\bibitem[5]{source5} Franz Baader and Tobias Nipkow. Term rewriting and all that. Cambridge university press, 1998.
\bibitem[6]{source6} Mouhebeddine Berrima, Narjes Ben Rajeb, and Véronique Cortier. Deciding knowledge in security protocols under some e-voting theories. RAIRO-Theoretical Informatics and Applications, 45(3):269-299, 2011. URL: https://doi.org/10.1051/ITA/2011119 .
\bibitem[7]{source7} Bruno Blanchet, Martín Abadi, and Cédric Fournet. Automated verification of selected equivalences for security protocols. The Journal of Logic and Algebraic Programming, 75(1):3-51, 2008. URL: https://doi.org/10.1016/J.JLAP.2007.06.002 .
\bibitem[8]{source8} Johannes Borgström. Static equivalence is harder than knowledge. Electronic Notes in Theoretical Computer Science, 154(3):45-57, 2006. URL: https://doi.org/10.1016/J.ENTCS.2006.05.006 .
\bibitem[9]{source9} Carter Bunch, Saraid Dwyer Satterfield, Serdar Erbatur, Andrew M Marshall, and Christophe Ringeissen. Knowledge problems in protocol analysis: Extending the notion of subterm convergent. arXiv preprint arXiv:2401.17226, 2024. URL: https://doi.org/10.48550/arXiv.2401.17226 .
\bibitem[10]{source10} Rohit Chadha, Vincent Cheval, Ştefan Ciobâcă, and Steve Kremer. Automated verification of equivalence properties of cryptographic protocols. ACM Transactions on Computational Logic (TOCL), 17(4):1-32, 2016.
\bibitem[11]{source11} Vincent Cheval. APTE: an algorithm for proving trace equivalence. In Erika Ábrahám and Klaus Havelund, editors, Tools and Algorithms for the Construction and Analysis of Systems - 20th International Conference, TACAS 2014. Proceedings, volume 8413 of Lecture Notes in Computer Science, pages 587-592. Springer, 2014. URL: https://doi.org/10.1007/978-3-642-54862-8\_50 .
\bibitem[12]{source12} Vincent Cheval, Steve Kremer, and Itsaka Rakotonirina. Deepsec: deciding equivalence properties in security protocols theory and practice. In 2018 IEEE symposium on security and privacy (SP), pages 529-546. IEEE, 2018. URL: https://doi.org/10.1109/SP.2018.00033 .
\bibitem[13]{source13} Ştefan Ciobâcă, Stéphanie Delaune, and Steve Kremer. Computing knowledge in security protocols under convergent equational theories. Journal of Automated Reasoning, 48(2):219-262, 2012. URL: https://doi.org/10.1007/S10817-010-9197-7 .
\bibitem[14]{source14} Katriel Cohn-Gordon, Cas Cremers, Luke Garratt, Jon Millican, and Kevin Milner. On ends-to-ends encryption: Asynchronous group messaging with strong security guarantees. In Proceedings of the 2018 ACM SIGSAC Conference on Computer and Communications Security, pages 1802-1819, 2018. URL: https://doi.org/10.1145/3243734.3243747 .
\bibitem[15]{source15} Véronique Cortier, Antoine Dallon, and Stéphanie Delaune. Sat-equiv: An efficient tool for equivalence properties. In 30th IEEE Computer Security Foundations Symposium, CSF 2017, pages 481-494. IEEE Computer Society, 2017. URL: https://doi.org/10.1109/CSF.2017.15 .
\bibitem[16]{source16} Véronique Cortier, Stéphanie Delaune, and Pascal Lafourcade. A survey of algebraic properties used in cryptographic protocols. Journal of Computer Security, 14(1):1-43, 2006. URL: https://doi.org/10.3233/JCS-2006-14101 .
\bibitem[17]{source17} Danny Dolev and Andrew Yao. On the security of public key protocols. IEEE Transactions on information theory, 29(2):198-208, 1983. URL: https://doi.org/10.1109/TIT.1983.1056650 .
\bibitem[18]{source18} Serdar Erbatur, Andrew M Marshall, Paliath Narendran, and Christophe Ringeissen. Deciding knowledge problems modulo classes of permutative theories. In International Symposium on Logic-Based Program Synthesis and Transformation, pages 47-63. Springer, 2024. URL: https://doi.org/10.1007/978-3-031-71294-4\_3 .
\bibitem[19]{source19} Serdar Erbatur, Andrew M. Marshall, Paliath Narendran, and Christophe Ringeissen. Knowledge Problems vs Unification and Matching: Dichotomy Results. In 10th International Conference on Formal Structures for Computation and Deduction (FSCD 2025), volume 337 of Leibniz International Proceedings in Informatics (LIPIcs), pages 18:1-18:17. Schloss Dagstuhl - Leibniz-Zentrum für Informatik, 2025. URL: https://doi.org/10.4230/LIPIcs.FSCD.2025.18 .
\bibitem[20]{source20} Serdar Erbatur, Andrew M Marshall, and Christophe Ringeissen. Notions of knowledge in combinations of theories sharing constructors. In International Conference on Automated Deduction, pages 60-76. Springer, 2017. URL: https://doi.org/10.1007/978-3-319-63046-5\_5 .
\bibitem[21]{source21} Serdar Erbatur, Andrew M Marshall, and Christophe Ringeissen. Computing knowledge in equational extensions of subterm convergent theories. Mathematical Structures in Computer Science, 30(6):683-709, 2020. URL: https://doi.org/10.1017/S0960129520000031 .
\bibitem[22]{source22} Paliath Narendran and Friedrich Otto. Single versus simultaneous equational unification and equational unification for variable-permuting theories. Journal of Automated Reasoning, 19(1):87-115, 1997. URL: https://doi.org/10.1023/A:1005764526878 .
\bibitem[23]{source23} Ky Nguyen. Formal verification of a messaging protocol. Work done under the supervision of Vincent Cheval and Véronique Cortier, 2019.
\bibitem[24]{source24} Saraid Dwyer Satterfield, Serdar Erbatur, Andrew M Marshall, and Christophe Ringeissen. Knowledge problems in security protocols: going beyond subterm convergent theories. In 8th International Conference on Formal Structures for Computation and Deduction (FSCD 2023), volume 260, pages 30-1, 2023.
\bibitem[25]{source25} Manfred Schmidt-Schauß. Unification in permutative equational theories is undecidable. Journal of symbolic computation, 8(4):415-421, 1989. URL: https://doi.org/10.1016/S0747-7171(89)80037-9 .
\bibitem[26]{source26} Alwen Tiu, Nam Nguyen, and Ross Horne. SPEC: an equivalence checker for security protocols. In Atsushi Igarashi, editor, Programming Languages and Systems - 14th Asian Symposium, APLAS 2016, Hanoi, Vietnam, November 21-23, 2016, Proceedings, volume 10017 of Lecture Notes in Computer Science, pages 87-95, 2016. URL: https://doi.org/10.1007/978-3-319-47958-3\_5 .
\end{thebibliography}
